\documentclass[10pt,a4paper]{amsart}
\usepackage[margin=19mm]{geometry}
\usepackage{amsmath,amssymb,mathtools}
\usepackage[scr=rsfs,cal=euler]{mathalfa}
\usepackage{xfrac}
\usepackage[unicode,hidelinks,bookmarks=false]{hyperref}
\newcommand{\ie}{i.e.\ }
\newcommand{\cf}{cf.\ }
\newcommand{\resp}{respectively\ }
\newcommand{\Subsection}{\subsection}
\providecommand{\nobreakdash}{\nobreak\hbox{-}}
\setlength{\emergencystretch}{2em}
\multlinegap0pt
\usepackage{bm}
\usepackage{bbm}
\usepackage{dsfont}
\usepackage{xspace}
\usepackage{cancel}
\usepackage{mathtools}
\usepackage{amsthm}
\makeatletter
\@ifundefined{definition}{\newtheorem{definition}{Definition}}{}
\@ifundefined{proposition}{\newtheorem{proposition}{Proposition}}{}
\@ifundefined{theorem}{\newtheorem{theorem}{Theorem}}{}
\makeatother
\title[Formalising Propositional Information via Implication Hyperg]{Formalising Propositional Information via Implication Hypergraphs}
\author{Vibhu Dalal}
\date{Source: arXiv:2502.00186v2 (CC BY 4.0)}
\begin{document}
\maketitle

\noindent\emph{Source and licence.} Formalising Propositional Information via Implication Hypergraphs. Version arXiv:2502.00186v2 (CC BY 4.0), distributed under the Creative Commons Attribution 4.0 International licence, as recorded in the arXiv licence metadata for this exact version and shown on the abstract page https://arxiv.org/abs/2502.00186v2. Full rights notice: licenses/arxiv-2502.00186-CC-BY-4.0.txt.\par\smallskip

\noindent\textit{Editorial scope.} Counted unit: the Theorem whose source label is \texttt{thm:example2} in the article (source0.tex lines 298--301, source label \texttt{thm:example2}), delivered together with the article's own complete original proof of it (source0.tex lines 303--328, one \texttt{proof} environment of 26 lines). No mathematics is written, repaired or re-derived: statement and proof are the article's own text line for line. Required context retained uncounted: def:prop\_info (lines 170--177); prop:well-defined (lines 205--212). Every definition, display and label of those lines is retained unchanged. Omitted from this package and not redistributed: 1--169, 178--204, 213--297, 302--302, 329--494. Nothing inside a retained excerpt is omitted or altered: every retained body is delivered exactly as the article's lines are written, so no omission receipt of any kind accompanies this package. Citations: the retained excerpts carry no citation command at all, so no bibliography entry of the article is transcribed and no reference is redistributed. Reference adaptation: none was needed and none was made; every reference inside the delivered unit resolves inside it, so no unresolved reference or citation is produced. Printed slips kept literal: no printed word, symbol, hypothesis or number of the counted statement or of its proof was altered. Layout and class: the article's own class is \texttt{article} with packages including geometry, amsmath, amssymb, amsthm, graphicx, enumitem, mathpazo, titlesec, hyperref, nicefrac, float, cite; the wrapper is the lane's pinned \texttt{amsart} editorial wrapper at 10pt on A4, which restates the theorem-like environments the retained text needs and adds the article's own macro definitions that the retained text needs (none), with extra wrapper packages (bm, bbm, dsfont, xspace, cancel, mathtools). The delivered numbering is the unit's own. Assets: no retained span contains a figure environment, a graphics-inclusion command, a PDF-page inclusion, a table or a TikZ picture; no image file is copied into this package, the package declares no asset dependency and no image file is redistributed with it. Licence: version arXiv:2502.00186v2 of this article is distributed under the Creative Commons Attribution 4.0 International licence, as recorded in the arXiv licence metadata for this exact version and shown on the abstract page https://arxiv.org/abs/2502.00186v2. A full rights notice is delivered at licenses/arxiv-2502.00186-CC-BY-4.0.txt. Characters: every retained excerpt is pure ASCII, so the delivered engine, XeLaTeX with its default Latin Modern text font, reports no missing character.
\begin{definition}[Adjacency Matrix]\label{def:prop_info}
Let \( H = (V, E) \) be an implication hypergraph. We denote its vertices by \( V = \{p_1,p_2,\ldots,p_n\} \). 
Let \( E_i = \{e \in E \mid p_i \in h(e)\} \). Then the adjacency matrix of \( H \), \( A \), is defined by:
\[
A_{ij} = \sum_{\substack{e \in E_i}} \frac{1}{|h(e)|} \mathbf{1}_{p_j \in t(e)}
\]
for \( 1 \leq i, j \leq n \).
\end{definition}

\begin{proposition}\label{prop:well-defined}
\normalfont
Let $H\in \mathcal{H}$, and $\nu, \epsilon>0$. Let $A$ denote the adjacency matrix of $H=(V,E)$, where $V=\{p_1,\ldots,p_n\}$. Let \( \mathbf{1} \) be the vector of $n$ ones, and \( l \) be a vector where the \( i \)-th entry is 1 if \( p_i \) is a leaf node, and 0 otherwise. If \( A - I \) is invertible, or equivalently if 1 is not an eigenvalue of \( A \), then \( \mathcal{I}_{\nu, \epsilon} = (\mathcal{I}_{\nu, \epsilon}(p_1), \ldots, \mathcal{I}_{\nu, \epsilon}(p_n)) \) is well-defined, and

\[
\mathcal{I}_{\nu, \epsilon} = (I - A)^{-1} \left( \epsilon A \mathbf{1} + \nu l \right).
\]
\end{proposition}

\begin{theorem}\label{thm:example2}
\normalfont
Let \( H=(V, E)\in \mathcal{H} \), where \( V = \{p_1, p_2, \ldots, p_n\} \). Let \( A \) denote its adjacency matrix. If for all \( k \in \{1, \ldots, n\} \), the diagonal of \( A^k \) is zero, i.e., \( \text{diag}(A^k) = \mathbf{0} \), then \( H \) is \emph{configured}.
\end{theorem}

\begin{proof}
The case \( n = 1 \) is straightforward, as the graph consists of a single node, which is necessarily a leaf node. Consequently, its information value is \( \nu > 0 \). Now, let us consider the case where \( n \geq 2 \).

We prove by contradiction that \( 1 \) cannot be an eigenvalue of \( A \). Assume, for the sake of contradiction, that \( 1 \) is an eigenvalue of \( A \). Since the eigenvalues of \( A^2 \) are the squares of the eigenvalues of \( A \), this would imply that \( \text{trace}(A^2) > 0 \). However, by our hypothesis, we have \( \text{diag}(A^2) = \mathbf{0} \), leading to a contradiction. Therefore, \( 1 \) cannot be an eigenvalue of \( A \). Consequently, by Proposition~\ref{prop:well-defined}, \( \mathcal{I}_{\nu, \epsilon} = (\mathcal{I}_{\nu, \epsilon}(p_1), \ldots, \mathcal{I}_{\nu, \epsilon}(p_n)) \) is well-defined.

Let $p_i \in V$. If $p_i \in L(H)$, then $I_{\nu, \epsilon}(p_i) = \nu > 0$. Suppose $p_i \notin L(H)$. To show that $I_{\nu, \epsilon}(p_i) > 0$, we proceed by contradiction. Suppose that $I_{\nu, \epsilon}(p_i) \leq 0$. By Definition~\ref{def:prop_info}:
\[
I_{\nu, \epsilon}(p_i)=\sum_{j=1}^{n} A_{ij} (\mathcal{I}_{\nu, \epsilon}(p_j) + \epsilon).
\]
For all \( i, j \in \{1, \ldots, n\} \), \( A_{ij} \geq 0 \). Let \( S \subset \{1, \ldots, n\} \) be the set of indices \( s \) such that \( A_{is} > 0 \) for all \( s \in S \). Since \( p_i \notin L(H) \), we have \( S \neq \emptyset \). Therefore, there exists \( i_1 \in S \) such that \( \mathcal{I}_{\nu,\epsilon}(p_{i_1}) \leq 0 \). Note that there exists a hyperedge \( e \in E \) such that \( i \in H(e) \) and \( i_1 \in T(e) \). By repeatedly applying this reasoning, we construct a sequence 
$Q = (p_{i_1}, p_{i_2}, \ldots)$ of vertices such that:
\begin{enumerate}
    \item For each pair of consecutive vertices \( p_{i_k}, p_{i_{k+1}} \) in \( Q \), there exists a hyperedge connecting them.
    \item \( \forall p_{i_k} \in Q \), \( \mathcal{I}_{\nu,\epsilon}(p_{i_k}) \leq 0 \).
    \item Q terminates at $p_{i_k}$ if and only if $p_{i_k}\in L(H)$.
\end{enumerate}

Since \( \operatorname{diag}(A^k) = \mathbf{0} \) for all \( k \in \{1, \ldots, n\} \), \( H \) is acyclic. Hence, the sequence \( Q \) must be finite. Let $ Q = (p_{i_1}, p_{i_2}, \ldots, p_{i_m}), $
where \( m < n \), and consequently \( p_{i_m} \in L(H) \). However, if \( p_{i_m} \in L(H) \), then 
$ \mathcal{I}_{\nu,\epsilon}(p_{i_m}) = \nu > 0, $
which contradicts our earlier assumption that \( \mathcal{I}_{\nu,\epsilon}(p_{i_k}) \leq 0 \) for all \( p_{i_k} \in Q \).

Thus, \( \mathcal{I}_{\nu,\epsilon}(p_i) > 0 \) holds for all \( p_i \notin L(H) \), completing the proof.


\end{proof}

\end{document}
